\chapter{Előzmények és irodalomkutatás}
\label{chap:irodalom}

A szimulációs rendszerek tervezése előtt alapos irodalomkutatást végeztem az ACM Digital Library adatbázisában \footnote{Az ACM (Association for Computing Machinery) Digital Library a világ egyik legjelentősebb informatikai tudományos adatbázisa, amely hozzáférést biztosít neves szakfolyóiratokhoz és konferenciakiadványokhoz: \url{https://dl.acm.org/}}, valamint megvizsgáltam a területen elérhető népszerű szoftveres megoldásokat és oktatási célú projekteket. A kutatás célja az volt, hogy feltérképezzem az ágensalapú modellezés (Agent-Based Modeling – ABM) aktuális módszereit és a döntéshozatali algoritmusok implementációs lehetőségeit.

% 3.1
%~~~~~~~~~~~~~~~~~~~~~~~~~~~~~~~~~~~~~~~~~~~~~~~~~~~~~~~~~~~~~~~~~~~~~~~~~~~~~~~~~~~~~~
\section{Irodalomkutatás és elméleti háttér}

Az ökoszisztémák modellezése és az ágensalapú szimulációk területén végzett kutatásaim során több olyan tudományos publikációt vizsgáltam meg, amelyek alapjaiban határozták meg a fejlesztés irányvonalát. Az alábbiakban ezeket a munkákat részletezem, kiemelve a saját megoldásomtól való eltéréseket.

Tiago és Reis munkájukban \cite{Tiago08} a koevolúció folyamatát vizsgálták, ahol a ragadozó és préda fajok egyidejűleg fejlődnek genetikus programozás és döntési fák segítségével. A cikk fő fókusza a stratégiai viselkedésformák kialakulása volt: a kutatók azt elemezték, hogy a fajok hogyan képesek komplex válaszokat adni egymás változó taktikáira. Bár a genetikus megközelítés náluk is központi szerepet játszik, az én megoldásom lényegesen eltér a viselkedés leírásának módjában. Míg ők döntési fákat generáltak evolúciós úton, én egy rögzített, de paramétereiben (genotípus) örökölhető és mutálható szabályrendszert alkottam meg. Ez lehetővé teszi a populációgenetikai adatok (példából: látótávolság, sebesség) sokkal finomabb, folytonos skálán történő követését és statisztikai elemzését, szemben a diszkrét döntési ágakkal.

Bearss és társai \cite{Michael22} a klasszikus, Wilensky-féle „Wolf-Sheep Predation” modellt ültették át modern környezetbe egy oktatási gyakorlat keretében. Tanulmányuk rávilágított arra, hogy az ágensalapú modellezés (ABM) hatékonysága nagyban függ a választott keretrendszer számítási kapacitásától és vizualizációs képességeitől. Az ő munkájuk elsősorban a modell validálására és az oktatási módszertanra fókuszált. Ezzel szemben az én fejlesztésem nem csupán egy meglévő modell újraimplementálása, hanem egy kiterjeszthető mérnöki keretrendszer létrehozása. Míg az említett munka megmarad a hagyományos 2D-s rácsalapú szemléletnél, az én szimulációm a Unity3D motor lehetőségeit kihasználva folytonos 3D-s térben, procedurálisan generált domborzaton zajlik, ami merőben más útvonalkeresési és interakciós logikát igényel.

A projekt szempontjából a legmeghatározóbb Gras és szerzőtársai publikációja \cite{Gras09} volt, amely a Fuzzy Cognitive Map (FCM) alkalmazását mutatja be egyén alapú ökoszisztéma-szimulációkban. A cikk központi gondolata, hogy a biológiai lények döntéshozatala nem bináris logikát követ, hanem belső állapotok (pl. félelem, éhség, fáradtság) és külső ingerek súlyozott eredője. Bár a viselkedésmodellezés alapelveit tőlük kölcsönöztem, a megvalósítás technikai háttere és a vizsgált módszerek köre jelentősen különbözik. Grasék kutatása egy dedikált, tudományos célú szoftverkörnyezetben készült, ahol a vizualizáció és a szoftverarchitektúra rugalmassága másodlagos volt. Ezzel szemben az én munkámban a hangsúly a szoftvertechnológiai modernizáción és a különböző döntési paradigmák összehasonlíthatóságán van.

A fejlesztés során egy olyan moduláris keretrendszert alakítottam ki, amely lehetővé teszi három különböző döntési mechanizmus párhuzamos vizsgálatát és összehasonlítását:

\begin{itemize}
    \item \textbf{Véges állapotgép (FSM):} A BSc szakaszban kidolgozott megoldás továbbvitele, ahol az ágensek viselkedését előre rögzített állapotátmeneti szabályok határozzák meg. Bár a módszer stabil, mérnöki szempontból viszont rugalmatlan, mivel nem teszi lehetővé az egyedi adaptációt és az evolúciós fejlődést.
    \item \textbf{Hasznossági alapú döntéshozatal (Utility AI):} Ebben a modellben az ágensek különböző súlyozott szempontok alapján számítják ki az egyes cselekvések várható hasznosságát. A súlyokat egy különálló, genotípust leíró osztály (\texttt{Genome}) tárolja, amely az öröklődés során mutálódni képes. Ez a megközelítés már lehetőséget ad az ökoszisztéma evolúciós fejlődésének megfigyelésére.
    \item \textbf{Fuzzy Cognitive Map (FCM):} A legkomplexebb alkalmazott módszer, amelyben a döntéshozatalt egy irányított, súlyozott gráf reprezentálja. Itt a belső állapotok és fogalmak egymásra gyakorolt hatása révén apró változások is finom, nem-lineáris viselkedésformákat eredményezhetnek. A rendszer tervezésekor külön figyelmet fordítottam ezen komplex gráfszerkezet hatékony tárolására és az eseményvezérelt architektúrába való integrálására.
\end{itemize}

A saját megoldásom újdonsága, hogy ezeket a mechanizmusokat egy egységes, eseményvezérelt Unity-architektúrába integráltam, kiegészítve real-time HTTP alapú naplózással. Ezáltal lehetőség nyílik arra, hogy a statisztikai elemzés során számszerűsíthető legyen az egyes "gondolkodásmódok" hatása a populáció túlélési rátájára és stabilitására.

A három architektúra közötti alapvető szoftvertechnológiai és elméleti különbségeket a \ref{table:decision_comparison}. táblázat foglalja össze részletesen:

\small
\begin{longtable}{| >{\RaggedRight}p{2.8cm} | >{\RaggedRight}p{3.8cm} | >{\RaggedRight}p{3.8cm} | >{\RaggedRight}p{3.8cm} |}
\caption{Az FSM, Utility AI és FCM döntési architektúrák összehasonlítása}
\label{table:decision_comparison} \\
\hline
\textbf{} & \textbf{Véges állapotgép \newline (FSM)} & \textbf{Hasznossági AI \newline (Utility)} & \textbf{Fuzzy Cognitive Map \newline (FCM)} \\ \hline

\endhead

\hline
\multicolumn{4}{r}{\textit{Folytatás a következő oldalon...}} \\
\endfoot

\hline
\endlastfoot

\textbf{Alapvető működési elv} & Diszkrét állapotok és determinisztikus átmenetek összessége. & Súlyozott matematikai függvények és hasznossági értékek maximumkeresése. & Koncepciók (belső állapotok és akciók) irányított, súlyozott gráfja. \\ \hline

\textbf{Döntéshozatali logika} & Feltételvezérelt (if-then), merev küszöbértékek alapján. & Folytonos, normalizált értékek lineáris kombinációja. & Iteratív, nem-lineáris állapotfrissítés szigmoid aktivációs függvénnyel. \\ \hline

\textbf{Viselkedési determinizmus} & Teljesen determinisztikus és statikus; azonos ingerre mindig azonos válasz. & Dinamikus: az egyedi genom súlyai miatt azonos állapot más döntést szülhet. & Emergens és kontextusfüggő: a gráfon belüli visszacsatolások egyedi mintákat adnak. \\ \hline

\textbf{Paraméterek reprezentációja} & \textit{Hardcoded} logikai szabályok és konstansok. & Elkülönített, örökölhető és mutálható leíró osztály: \texttt{UtilityGenome}. & Súlymatrixtól és topológiától függő egyedi tömb: \texttt{FCMGenome}. \\ \hline

\textbf{Evolúciós alkalmasság} & Nem alkalmas; a struktúra futásidőben nem alakítható vagy finomhangolható. & Kiválóan alkalmas; a súlyok folytonos skálán való mutációja stabilan követhető. & Összetett; a súlyok mutációja mellett a topológia stabilitása nehezen garantálható. \\ \hline

\textbf{Szoftver-architektúrális előnyök} & Rendkívül alacsony számítási igény, egyszerű és gyors implementáció. & Moduláris, könnyen skálázható, az ágensek szintjén jól elkülöníthető. & Komplex, nem-lineáris visszacsatolások és finom viselkedésformák modellezése. \\ \hline

\textbf{Mérnöki korlátok és hátrányok} & Merev, skálázhatatlan; komplex rendszernél „állapotrobbanás” lép fel. & Nem veszi figyelembe az egyes döntési tényezők egymásra gyakorolt hatásait. & Magasabb számítási és memóriaigény, összetett debugolás és finomhangolás. \\ \hline
\end{longtable}
\normalsize

% 3.2
%~~~~~~~~~~~~~~~~~~~~~~~~~~~~~~~~~~~~~~~~~~~~~~~~~~~~~~~~~~~~~~~~~~~~~~~~~~~~~~~~~~~~~~
\section{Hasonló alkotások és inspirációk}

A szoftverfejlesztői közösségben, különösen a Unity3D motor köré csoportosulva, több figyelemre méltó projekt is foglalkozik ökológiai szimulációval. Ezek az alkotások jelentős inspirációt nyújtottak a fejlesztés során, ugyanakkor rávilágítottak olyan technikai korlátokra is, amelyeket jelen munka igyekszik feloldani.

\subsection{Sebastian Lague: Simulating Ecosystems}

A projekt elsődleges kiindulópontját Sebastian Lague munkássága \cite{Sebastian19} jelentette. Az általa bemutatott szimuláció sikeresen demonstrálta az ágensalapú modellezés alapvető törvényszerűségeit egy letisztult vizuális környezetben. 

Bár jelen dolgozat alapgondolatai sok ponton hasonlítanak Lague megoldásához, szoftvertechnológiai és környezeti szempontból jelentős eltérések mutatkoznak. Míg Lague implementációja elsősorban a koncepció közérthető bemutatására és az egyszerűségre fókuszált, jelen rendszer egy robusztusabb, ipari sztenderdekhez közelebb álló mérnöki megközelítést alkalmaz. 

Környezeti szempontból Lague munkája egy sík, határolt 3D-s térben zajlik, ezzel szemben az én rendszeremben az ágensek egy komplex, procedurálisan generált 3D-s domborzaton mozognak. Ez jelentősen megnöveli a navigációs algoritmusok (NavMesh) és a térbeli interakciók számítási igényét.

Szoftverarchitekturális szempontból a legfontosabb különbség az ágensek közötti kommunikáció és a környezetérzékelés megvalósítása. Lague szimulációja erősen támaszkodik a folytonos lekérdezésre (polling): az ágensek a Unity \texttt{Update()} és \texttt{FixedUpdate()} ciklusaiban, folyamatosan futó távolságmérésekkel és környezeti lekérdezésekkel határozzák meg a környezetük állapotát, ami magas ágensszám esetén jelentős teljesítménycsökkenéshez vezet. Jelen fejlesztés során ezt egy eseményvezérelt működéssel (C\# eventek és delegátumok használatával) váltottam ki, ahol az ágensek állapota és döntési logikája csak releváns események (például táplálék generálódása vagy egy másik ágens pusztulása) hatására frissül. Ez a megközelítés drasztikusan javítja az erőforrás-kezelést és a skálázhatóságot.

Emellett a szimuláció megfigyelhetősége is szintet lépett: a statisztikai adatok gyűjtése nem egyszerű lokális fájlba írással történik, hanem egy professzionális idősoros adatbázisba (InfluxDB) továbbítódik HTTP protokollon keresztül. Ez lehetővé tette a Grafana alapú dashboardok használatát, ahol a populációdinamikai mutatók valós időben, interaktív módon elemezhetők, messze túlmutatva a hivatkozott munka vizualizációs lehetőségein.

\subsection{VAV: Evolúciós devlog}

VAV fejlesztése \cite{VAV23} Lague alapjaira építve mutat be egy továbbfejlesztett modellt. Ebben a megoldásban kiemelkedő a morfológiai jellemzők — például az állatok nyakhosszának — öröklődése és környezethez való adaptációja. 

Jelen munka szempontjából ezen projekt tanulsága az volt, hogy a vizuális jegyek és a funkcionális tulajdonságok közötti kapcsolat (például a testfelépítés és az elért táplálék viszonya) jelentősen növeli a szimuláció biológiai hűségét. Noha a jelenlegi fázisban a hangsúly inkább a populációgenetikai adatok statisztikai elemzésén és a szoftverarchitektúrán van, a VAV által bemutatott morfológiai diverzitás a későbbi továbbfejlesztési irányok egyik alapkövét jelenti.

\subsection{Fun Master Ed}

Fun Master Ed munkái számomra külön kiemelkednek az ökoszisztéma-szimulációk közül azáltal, hogy nem ragaszkodnak a konvencionális, valósághű ábrázolásmódhoz, hanem absztrakt modellek segítségével vizsgálnak komplex biológiai és pszichológiai jelenségeket.

Első munkájában \cite{Fun20} a hangsúly a szociális és érzelmi alapú viselkedésmodellezésen van. A szimuláció olyan változókat vezet be, mint a „boldogság”, valamint a passzív és agresszív viselkedési attitűdök, vizsgálva ezek hatását a populáció stabilitására és a ragadozó-préda interakciókra. Bár az én fejlesztésem elsősorban a túlélési és reprodukciós stratégiákra fókuszál az érzelmi alapú állapotok helyett, ez a megközelítés rávilágított arra, hogy a döntési folyamatokat érdemes belső, rejtett paraméterekkel is befolyásolni, ami némiképp rokonítható a nálam alkalmazott Fuzzy Cognitive Map belső állapotaival.

A szerző későbbi projektje \cite{Fun22} egy másik kritikus irányba, a morfológiai adaptáció felé mutat. Ebben a szimulációban az ágensek testfelépítése (például a lábak száma vagy a nyak magassága) közvetlenül a környezeti változásokra adott válaszként, az evolúció útján alakul ki. Jó példa erre a zsiráfok evolúciója: a magasabban elérhető táplálékforrás szelekciós nyomást gyakorolt a nyakhossz növekedésére. 

Bár a jelen szimuláció ragadozó-préda modellje (rókák és nyulak) konkrét biológiai analógiákra épít, a hivatkozott munkák rávilágítottak az absztrakt modellezés előnyeire és a környezethez való fizikai/mentális alkalmazkodás fontosságára. Ezen inspirációk hatására alakítottam ki a döntéshozatali logikát moduláris módon: a rendszeremben az ágensek „gondolkodásmódja” (FSM, Utility vagy FCM) és a hozzájuk tartozó genetikai paraméterek (\texttt{Genome}) elválnak a fizikai reprezentációtól. Ez a struktúra biztosítja, hogy a jövőben ne csak előre rögzített fajok, hanem dinamikusan változó viselkedésmintájú és testfelépítésű ágensek is benépesíthessék a szimulációs teret.

\subsection{Primer}

A technikai implementációk mellett Primer youtube csatorna videósorozata \cite{Primer18} nyújtott jelentős elméleti hátteret a szimuláció biológiai szabályrendszerének kialakításához. Primer munkásságának fő fókusza nem a vizuális komplexitás, hanem az evolúciós folyamatok matematikai és játékelméleti modellezése. Szimulációiban az ágensek (úgynevezett „Blob”-ok) egyszerű szabályok szerint mozognak, azonban a mutációk révén kialakuló tulajdonságok — mint a sebesség vagy a méret — közvetlen hatással vannak az energiafelhasználásra és a túlélési esélyekre.

Az ő megközelítése segített validálni a saját modellmben alkalmazott „trade-off” mechanizmusokat. Primer rávilágított arra, hogy egy sikeres szimulációban minden előnyös tulajdonságnak (pl. nagyobb sebesség) rendelkeznie kell egy hátránnyal is (pl. magasabb alapanyagcsere-igény), különben a populáció instabillá válik. Bár Primer megoldásai absztrakt, rácsalapú vagy egyszerűsített 2D-s tereken futnak, a tőle átvett elméleti összefüggéseket én egy komplex, 3D-s Unity környezetbe ültettem át. Ez a váltás lehetővé tette, hogy az ő statisztikai alapú következtetéseit egy sokkal dinamikusabb, fizikai alapú interakciókkal (pl. tényleges ütközések, domborzati viszonyok) rendelkező rendszerben vizsgáljam tovább.

% 3.3
%~~~~~~~~~~~~~~~~~~~~~~~~~~~~~~~~~~~~~~~~~~~~~~~~~~~~~~~~~~~~~~~~~~~~~~~~~~~~~~~~~~~~~~
\section{Következtetések és tervezési irányelvek}

Az irodalomkutatás, valamint a létező szoftveres megoldások és devlogok elemzése alapján egyértelművé vált, hogy a biológiailag hiteles és szoftvertechnológiailag is fenntartható ökoszisztéma-szimuláció egyedi mérnöki megközelítést igényel. A vizsgált munkákból levont tanulságokat az alábbi konkrét tervezési döntések formájában integrálom a saját rendszerembe:

\begin{enumerate}
    \item \textbf{A döntési mechanizmusok hibridizációja:} Gras és szerzőtársai \cite{Gras09} munkája igazolta a Fuzzy Cognitive Map (FCM) fölényét a merev struktúrákkal szemben, míg Tiago és Reis \cite{Tiago08} a diszkrét logikai döntések evolúcióját vizsgálták. E két megközelítést szintetizálva a saját rendszeremben nem egyetlen modellt alkalmazok, hanem egy olyan összehasonlító keretrendszert tervezek, amely az FSM, a folytonos genetikai térrel rendelkező Utility AI és az FCM modellek teljesítményét azonos környezetben, számszerűsíthető módon veti össze.
    
    \item \textbf{A polling kiváltása eseményvezérelt architektúrával:} Sebastian Lague \cite{Sebastian19} megoldásának kritikai elemzése rávilágított arra, hogy a játékmotorok hagyományos \texttt{Update()}-alapú polling mechanizmusai komoly teljesítménybeli korlátot jelentenek magas ágensszám mellett. A skálázhatóság érdekében a saját rendszerem tervezésekor alapvető elvárás a teljes körű eseményvezéreltség (C\# event-driven megközelítés), ahol az ágensek érzékelési és döntési ciklusai csak indokolt környezeti változások hatására futnak le.
    
    \item \textbf{Komplex térbeli és morfológiai modellezés:} Szemben a Bearss \cite{Michael22} által bemutatott hagyományos, kétdimenziós rácsalapú modellekkel, jelen munka a Unity3D motor lehetőségeit kihasználva folytonos 3D-s térben, procedurálisan generált domborzaton valósul meg. VAV \cite{VAV23} és Fun Master Ed \cite{Fun22} munkáiból kiindulva a fizikai környezet (domborzat) és a funkcionális tulajdonságok kapcsolatát úgy tervezem meg, hogy a terepviszonyok közvetlen szelekciós nyomást gyakoroljanak az olyan öröklődő tulajdonságokra, mint a sebesség és az energiafelhasználás.
    
    \item \textbf{Biológiai egyensúly és trade-off mechanizmusok:} Primer \cite{Primer18} szimulációi matematikai pontossággal igazolták, hogy a populáció stabilitásához elengedhetetlen az egymással versengő tulajdonságok (trade-offok) bevezetése. A \texttt{Genome} osztály tervezésekor ezt az elvet alkalmazom: a nagyobb mozgási sebesség vagy a kiterjedtebb látómező magasabb alapanyagcsere- és energiaköltséggel fog járni, megakadályozva az irreális „szuper-ágensek” kialakulását és a szimuláció korai összeomlását.
    
    \item \textbf{Telemetria és tudományos megfigyelhetőség:} Míg a vizsgált YouTube-projektek többsége \cite{Sebastian19, VAV23, Fun20} beért a vizuális reprezentációval vagy egyszerű lokális naplózással, egy MSc szintű kutatás megköveteli az adatok tudományos igényű validálását. Emiatt a rendszer alapvető részévé teszek egy HTTP-alapú külső telemetriai csatlakozót, amely az adatokat InfluxDB idősoros adatbázisba rendezi, és Grafana dashboardokon keresztül teszi mérhetővé a Lotka–Volterra-féle populáció-oszcillációkat.
\end{enumerate}